%
% schluss.tex -- Rueckblick und Ausblick
%
% (c) 2026 Baris Catan, OST Ostschweizer Fachhochschule
%
% !TEX root = ../../buch.tex
% !TEX encoding = UTF-8
%
\section{Eine folgenreiche Antwort
\label{elliptisch:section:schluss}}
\kopfrechts{Eine folgenreiche Antwort}

Die harmlose Frage nach dem Umfang der Ellipse hat mehrere grosse
Gebiete der Mathematik hervorgebracht.
Die Frage nach der elementaren Lösbarkeit ist zunächst in der
Analysis formuliert und mit dem Satz von Liouville
\eqref{elliptisch:equation:liouville} in der Algebra entschieden
worden.

Untersucht haben wir die Grenze an Integralen des Typs
$R\bigl(x,\sqrt{p(x)}\bigr)$ und daran, wo ihre elementare
Lösbarkeit endet.
Über den Reduktionssatz~\ref{elliptisch:satz:reduktion} sind daraus
die drei vollständigen Integrale $K$, $E$ und $\Pi$ entstanden.
Lässt man die obere Grenze variieren, so werden daraus die
unvollständigen Integrale, die historisch unhandlich blieben.
Den Durchbruch brachte erst ihre geschickte Umkehrung, aus der die
jacobischen elliptischen Funktionen $\operatorname{sn}$,
$\operatorname{cn}$ und $\operatorname{dn}$ hervorgegangen sind.
Sie erfüllen Identitäten und Additionstheoreme wie die bekannten
Kreisfunktionen, haben einfache Ableitungen und sind Lösungen der
Differentialgleichungen aus
Abschnitt~\ref{elliptisch:section:anwendungen}.
Nur kurz eingegangen sind wir auf ihre doppelte Periodizität, die
sich erst in der komplexen Ebene zeigt, wo diese Funktionen auch
vollständig zu verstehen sind.
Von dort aus haben die Theorie der riemannschen Flächen, die
Modulformen und die Arithmetik der elliptischen Kurven ihren
Aufschwung genommen, deren Anwendungen heute bis in die
Kryptographie reichen.
\index{riemannsche Flache@riemannsche Fläche}%
\index{Modulform}%

Dieses Kapitel deckt nur einen Bruchteil dessen ab, was aus dieser
Richtung gewachsen ist, von den Anwendungen bis zu ganzen neuen
Zweigen der Mathematik.
Was am Anfang wie eine Grenze der Analysis aussah, erwies sich als
Eingang zu einer eigenen Theorie.
